\section{The relation with the full cubic toroid}
Let $Q$ be the signed permutation group on three coordinates and
$D=\{v\in\Z^3:v_1+v_2+v_3\equiv0\pmod2\}$. In the notation of the
reviewed cubic-family proof, the hypertope group is $E_s=E/sD$, where
\[
 E=\{(q,v)\in Q\ltimes\Z^3:\textstyle\sum v_i\equiv\epsilon(q)\pmod2\},
 \qquad |E_s|=48s^3.
\]
The introductory sentence on page5 of \cite{source} identifies this group with
the full group of the cubic toroid $\{4,3,4\}_{(s,s,0)}$. The earlier
account \cite[pp.~14--15]{earlier} correctly describes an index-two subgroup.
Here is the exact change of generators and coordinates.

The standard Euclidean realization of $[4,3,4]$ is generated by
\begin{align*}
 R_0(x)&=(1-x_1,x_2,x_3),&R_1(x)&=(x_2,x_1,x_3),\\
 R_2(x)&=(x_1,x_3,x_2),&R_3(x)&=(x_1,x_2,-x_3).
\end{align*}
Its affine image is $W=Q\ltimes\Z^3$: the last three generators give
$Q$, and multiplication of $R_0$ by the corresponding linear reflection
gives translation by $e_1$. Conjugating gives all coordinate translations.
The faithful Euclidean Coxeter realization is a retained classical fact.
Direct multiplication gives
\[
 (R_0R_1R_2R_3R_2)^2=(I,e_1+e_2).
\]
Its $Q$-orbit generates $D$. Thus the additional relation with exponent
$2s$ used in \cite{earlier} imposes precisely $sD$, and the full toroid group is
$W/sD$, of order $48[\Z^3:sD]=96s^3$.

Put $p=(1/2,1/2,1/2)$ and $T(x)=\operatorname{diag}(1,-1,-1)(x-p)$.
Conjugating the four displayed ambient words by $T$ gives
\begin{align*}
 T(R_0R_1R_0)T^{-1}(x)&=(x_2,x_1,x_3),\\
 TR_2T^{-1}(x)&=(x_1,x_3,x_2),\\
 T(R_1R_3)T^{-1}(x)&=(-x_2,-x_1,1-x_3),\\
 TR_1T^{-1}(x)&=(-x_2,-x_1,x_3).
\end{align*}
These are exactly the four generators of $E$ in source-label order.
The linear part of $T$ preserves $D$, so this conjugacy descends to the
period quotients. Consequently the displayed subgroup of $W/sD$ is
isomorphic to $E/sD$ and has index two. In \cite{earlier} the Petrie generators
are $\alpha_0=R_0$, $\alpha_1=R_1$, $\alpha_2=R_2$ and
$\alpha_3=R_1R_3$; its four subgroup generators are the same set,
ordered here as $(\alpha_1^{\alpha_0},\alpha_2,\alpha_3,\alpha_1)$ to match \cite{source}.

The family still arises from the stated toroid. The correction is that
its group is the displayed index-two subgroup, rather than the full
toroid group. This changes no principal result of Section3 of \cite{source}.
It does not claim a new family or a resolution of an open question.
